%%%PAGE 139 rot=0 legibility=good summary=剪贴印刷题两道：毫安表改装与校准(含手写填空、双极限计算、判断理由)；星球弹簧a-x图像题(含手写批注)及印刷解析
%%%BLOCK id=139-01 ch=10 sec=电表改装·毫安表改装与校准 kind=problem exp=1 alt=none cont=none y=0.03-0.50
\begin{problem}[剪贴印刷题·9分]
图(a)为某同学改装和\fbox{校准}毫安表的电路图，其中虚线框内是毫安表的改装电路。% 原稿“校准”二字被手绘圈出
%FIG p139_a 0.05 0.048 0.31 0.146
\notefig[0.4]{figures/p139_a.png}
%FIG p139_b 0.33 0.07 0.49 0.16
\notefig[0.3]{figures/p139_b.png}

(1)已知毫安表表头的内阻为\fbox{$100\,\Omega$}，满偏电流为\fbox{$1\,\mathrm{mA}$}；$R_1$ 和 $R_2$ 为阻值固定的电阻。若使用 $a$ 和 $b$ 两个接线柱，电表量程为\fbox{$3\,\mathrm{mA}$}；若使用 $a$ 和 $c$ 两个接线柱，电表量程为\fbox{$10\,\mathrm{mA}$}。由题给条件和数据，可以求出 $R_1=$ \underline{\hspace{0.6em}\struck{35}15\hspace{0.6em}}\,$\Omega$，$R_2=$ \underline{\hspace{0.6em}35\hspace{0.6em}}\,$\Omega$。% 方框与空格里的数字均为手写；R1 空处原稿“35”似被划去后改写 15 (CHECK)

(2)现用一量程为 $3\,\mathrm{mA}$、内阻为 $150\,\Omega$ 的标准电流表\textcircled{A}对改装电表的 $3\,\mathrm{mA}$ 挡进行校准，校准时需选取的刻度为 $0.5$、$1.0$、$1.5$、$2.0$、$2.5$、$3.0\,\mathrm{mA}$。电池的电动势为 $1.5\,\mathrm{V}$，内阻忽略不计；定值电阻 $R_0$ 有两种规格，阻值分别为 $300\,\Omega$ 和 $1\,000\,\Omega$；滑动变阻器 $R$ 有两种规格，最大阻值分别为 $750\,\Omega$ 和 $3\,000\,\Omega$。则 $R_0$ 应选用阻值为 \underline{\hspace{0.4em}\struck{1000}300\hspace{0.4em}}\,$\Omega$ 的电阻，$R$ 应选用最大阻值为 \underline{\hspace{0.6em}3000\hspace{0.6em}}\,$\Omega$ 的滑动变阻器。% 原稿：(2)左侧有手绘括号“(”，“3 000 Ω。”后有“)”将条件括起；R0 空处为涂改，似为 1000 被划去后改写 300 (CHECK)

(3)若电阻 $R_1$ 和 $R_2$ 中有一个因损坏而阻值变为无穷大，利用图(b)的电路可以判断出损坏的电阻。图(b)中的 $R'$ 为保护电阻，虚线框内未画出的电路即为图(a)虚线框内的电路。则图中的 $d$ 点应和接线柱 \underline{\hspace{0.5em}C\hspace{0.5em}}(填“b”或“c”)相连。% 空格内手写 C

\begin{solution}
\textbf{双极限}% 手写，位于(2)问末尾右侧

最大刻度 $3\,\mathrm{mA}$\quad $\dfrac{E}{I}=500\,\Omega$\quad $R_A=150\,\Omega$\quad $\therefore R_0$ 用 $300\,\Omega$

最小刻度 $0.5\,\mathrm{mA}$\quad $\dfrac{E}{I}=3000\,\Omega$\quad 选 $3000\,\Omega$的

判断理由：闭合开关，若电表指针偏转，则损坏的电阻是 $R_1$；若电表指针不动，则\unclear{损}坏的是 $R_2$% CHECK: “损”字原稿形似“阻”

（(3)题“则图中的 $d$ 点…”一行左侧页边另有手写旁注：\unclear{双限}）
\end{solution}
\end{problem}
%%%END
%%%BLOCK id=139-02 ch=06 sec=弹簧模型·a-x图像 kind=problem alt=05,03 cont=none y=0.44-0.95
\begin{problem}[剪贴印刷题]
在星球 $M$ 上将一轻弹簧竖直固定在水平桌面上，把物体 $P$ 轻放在弹簧上端，$P$ 由静止向下运动，物体的加速度 $a$ 与弹簧的压缩量 $x$ 间的关系如图中实线所示。在另一星球 $N$ 上用完全相同的弹簧，改用物体 $Q$ 完成同样的过程，其 $a$—$x$ 关系如图中虚线所示。假设两星球均为质量均匀分布的球体。\underline{已知星球 $M$ 的半径是星球 $N$ 的 $3$ 倍，则}% 原稿此处有手画下划线
%FIG p139_c 0.31 0.455 0.57 0.637
\notefig[0.4]{figures/p139_c.png}

A. $M$ 与 $N$ 的密度相等

B. $Q$ 的质量是 $P$ 的 $3$ 倍\quad $M_Q=6M_P$% 手写（大写 M）

C. $Q$ 下落过程中的最大动能是 $P$ 的 $4$ 倍

D. $Q$ 下落过程中弹簧的最大压缩量是 $P$ 的 $4$ 倍\quad 对称性 \unclear{$\checkmark$}% 手写

\textbf{手写批注：}A\,\struck{B}\,C% CHECK: 大字手写 A B C，B 上有一竖线疑为划去（标准答案 AC）

$E_{pQ}=\dfrac{2a_0x_0}{2}$，\quad $E_{pP}=\dfrac{3a_0x_0}{2}$% CHECK: 下标原稿形似 pQ、pP，可能是 kQ、kP；图中纵轴 a 前被手写加了 m；图上另有一条红色手绘曲线

面积是能量

$E=max$

\begin{solution}
\textbf{21. AC}\quad 【命题意图】本题考查万有引力定律、牛顿第二定律和功能关系，意在考查考生的推理能力和分析综合能力。

【解题思路】设 $P$、$Q$ 的质量分别为 $m_P$、$m_Q$；$M$、$N$ 的质量分别为 $M_1$、$M_2$，半径分别为 $R_1$、$R_2$，密度分别为 $\rho_1$、$\rho_2$；$M$、$N$ 表面的重力加速度分别为 $g_1$、$g_2$。在星球 $M$ 上，弹簧压缩量为 $0$ 时有 $m_Pg_1=3m_Pa_0$，所以 $g_1=3a_0=G\dfrac{M_1}{R_1^2}$，密度 $\rho_1=\dfrac{M_1}{\frac{4}{3}\pi R_1^3}=\dfrac{9a_0}{4\pi GR_1}$；在星球 $N$ 上，弹簧压缩量为 $0$ 时有 $m_Qg_2=m_Qa_0$，所以 $g_2=a_0=G\dfrac{M_2}{R_2^2}$，密度 $\rho_2=\dfrac{M_2}{\frac{4}{3}\pi R_2^3}=\dfrac{3a_0}{4\pi GR_2}$；因为 $R_1=3R_2$，所以有 $\rho_1=\rho_2$，选项 A 正确。当物体的加速度为 $0$ 时有 $m_Pg_1=3m_Pa_0=kx_0$，$m_Qg_2=m_Qa_0=2kx_0$，解得 $m_Q=6m_P$，选项 B 错误。根据 $a$—$x$ 图线与坐标轴围成图形的面积和质量的乘积表示合外力做的功可知，$E_{\mathrm{km}P}=\dfrac{3}{2}m_Pa_0x_0$，$E_{\mathrm{km}Q}=m_Qa_0x_0$，所以 $E_{\mathrm{km}Q}=4E_{\mathrm{km}P}$，选项 C 正确。根据运动的对称性可知，$Q$ 下落时弹簧的最大压缩量为 $4x_0$，$P$ 下落时弹簧的最大压缩量为 $2x_0$，选项 D 错误。
\end{solution}
\end{problem}
%%%END
%%%PAGEEND 139
